% Auto-generated draft section (Stage 1 / P0): Golden Record formalism as full ER pipeline
% Path: 04_library/sources/research/2026/researchiq_latex_draft/sections/04_golden_record_formalism.tex

\section{Golden Record як повний конвеєр Entity Resolution (ER)}
\label{sec:golden-record-formalism}

У цьому розділі формалізовано побудову \emph{Golden Record} як відтворюваний наскрізний ER-конвеєр, що охоплює генерацію кандидатів (blocking), парне скоринг-оцінювання (record linkage), кластеризацію, злиття атрибутів з явним правилом розв’язання нічиїх (tie-breaking) та фіксацію провенансу.

\paragraph{Термінологічна примітка.}
Далі використовуємо такі відповідники: \textit{blocking} — \textbf{блокування} (генерація кандидатів), \textit{record linkage} — \textbf{зв’язування записів}, \textit{review} — \textbf{ручна перевірка}, \textit{tie-breaking} — \textbf{правило розв’язання нічиїх}.

\subsection{Постановка задачі та нотація}
\label{sec:gr-notation}

Нехай $\mathcal{S}$ --- скінченна множина джерел даних (наприклад, ORCID, OpenAlex, Crossref, локальна CRIS). Кожне джерело $s\in\mathcal{S}$ надає мультимножину сирих записів $R_s = \{r^{(s)}_1,\dots,r^{(s)}_{n_s}\}$. Сирий запис $r$ описує сутність реального світу типу $\tau\in\{\textsf{Publication},\textsf{Person},\textsf{Organization}\}$.

Кожен запис задається кортежем
\begin{equation}
  r = \big( \tau,\; \mathrm{id}_s,\; x,\; t,\; \pi \big),
\end{equation}
де $\mathrm{id}_s$ --- ідентифікатор запису, специфічний для джерела, $x$ --- (можливо розріджене) відображення атрибутів над множиною атрибутів $\mathcal{A}_\tau$, $t$ --- часовий штамп інжестингу, а $\pi$ --- необов’язковий нативний для джерела блок провенансу (наприклад, JSON).

\paragraph{Мета.}
За даними $R=\biguplus_{s\in\mathcal{S}} R_s$ побудувати:
(i) розбиття (або кластеризацію) $\mathcal{C}=\{c_1,\dots,c_K\}$ над (підмножиною) $R$ так, щоб кожен кластер відповідав одній базовій сутності, та
(ii) злитий канонічний запис (Golden Record) для кожного кластера, позначений $g(c_k)$, разом із провенансом на рівні атрибутів $\mathrm{prov}(g(c_k))$.

Ми розрізняємо два рівні:
\begin{itemize}
  \item \textbf{Зв’язування записів / матчинг (record linkage / matching):} оцінити парний бал відповідності $m(r_i,r_j)\in[0,1]$.
  \item \textbf{Кластеризація сутностей:} вивести кластери з вагового графа кандидатів, використовуючи $m(\cdot,\cdot)$.
\end{itemize}

\subsection{Канонікалізація та нормалізація}
\label{sec:gr-normalization}

Визначимо детерміновану функцію нормалізації $\nu_\tau$, яка відображає сирі значення атрибутів у канонічні форми, що використовуються для блокування та скорингу:
\begin{equation}
  \tilde{x} = \nu_\tau(x;\,\Theta_{\nu,\tau}),
\end{equation}
де $\Theta_{\nu,\tau}$ --- гіперпараметри нормалізації.

Для $\tau=\textsf{Publication}$ типовими компонентами є:
\begin{itemize}
  \item Нормалізація DOI: приведення до нижнього регістру, видалення префіксів (\texttt{https://doi.org/}, \texttt{dx.doi.org/}), валідація регулярним виразом.
  \item Нормалізація назви (title): нормалізація Unicode (NFKC), casefold, вилучення пунктуації, зведення пробілів, опційне вилучення стоп-слів.
  \item Нормалізація видання/журналу (venue): розгортання поширених скорочень (опційно), casefold, вилучення пунктуації.
  \item Нормалізація імен авторів: casefold, вилучення діакритиків, евристики перестановки токенів.
\end{itemize}

\paragraph{Відтворюваність.}
Усі кроки нормалізації мають бути версіоновані (наприклад, $\mathrm{ver}(\nu_\tau)$) та виконуватися детерміновано за заданих $\Theta_{\nu,\tau}$.

\subsection{Блокування (генерація кандидатів)}
\label{sec:gr-blocking}

Нехай $\mathcal{B}$ позначає множину функцій блокування. Функція блокування $b\in\mathcal{B}$ відображає запис у один або кілька ключів блокування:
\begin{equation}
  b(r) \rightarrow \{k_1,\dots,k_{\ell}\}, \quad k\in\Sigma^*.
\end{equation}
Кандидатні пари генеруються лише всередині блоків:
\begin{equation}
  \mathcal{P} = \bigcup_{b\in\mathcal{B}} \bigcup_{k} \{(r_i,r_j): r_i\neq r_j,\; k\in b(r_i)\cap b(r_j)\}.
\end{equation}

\paragraph{Політика блокування для публікацій (детермінований базовий варіант).}
Рекомендуємо двоступеневу схему блокування:
\begin{enumerate}
  \item \textbf{Точне блокування за DOI} (висока precision, висока recall за наявності DOI):
  \begin{equation}
    b_{\mathrm{doi}}(r)=\begin{cases}
      \{\mathrm{doi}(r)\} & \text{якщо } \mathrm{doi}(r) \text{ існує і є валідним}\\
      \emptyset & \text{інакше.}
    \end{cases}
  \end{equation}
  \item \textbf{Блокування за Title--Year} для відсутнього/невалідного DOI:
  \begin{equation}
    b_{\mathrm{ty}}(r)=\{\mathrm{minhash}_q(\mathrm{title}(r)) \;\|\|\; q\in\mathcal{Q}\}\times\{\mathrm{year}(r)+\delta: \delta\in\{-1,0,+1\}\},
  \end{equation}
  де $\mathcal{Q}$ --- множина розмірів $q$-грам (наприклад, $\{3,4\}$), а $\mathrm{minhash}_q$ --- детермінована схема хешування.
\end{enumerate}

\paragraph{Налаштовувані параметри (потрібна емпірична калібровка).}
Блокування задає компроміс між recall і обчислювальною вартістю. На розміченому development-наборі (або через ручний аудит) необхідно налаштувати:
(i) $\mathcal{Q}$ та конфігурацію хешування,
(ii) вікно за роком $\Delta_y$ (типово $\pm 1$),
(iii) опційне top-$K$ вилучення, якщо для генерації кандидатів використовується лексичний (BM25) або embedding-пошук.

\subsection{Парний скоринг (зв’язування записів)}
\label{sec:gr-pairwise}

Для кандидатної пари $(r_i,r_j)\in\mathcal{P}$ визначимо вектор ознак $\phi(r_i,r_j)\in\mathbb{R}^d$ на нормалізованих атрибутах (наприклад, рівність DOI, подібність назв, перетин авторів). Бал відповідності задається як:
\begin{equation}
  m(r_i,r_j)=\sigma\Big(\beta_0 + \sum_{p=1}^{d} \beta_p\,\phi_p(r_i,r_j)\Big),
  \label{eq:match-score}
\end{equation}
де $\sigma(z) = (1+e^{-z})^{-1}$.

\paragraph{Два операційні режими.}
\begin{itemize}
  \item \textbf{Правиловий скоринг (повністю детермінований):} обрати фіксовані ваги $\beta$.
  \item \textbf{Навчений скоринг (рекомендовано за наявності gold-розмітки):} оцінити $\beta$ за допомогою логістичної регресії (або подібної моделі) на розмічених парах.
\end{itemize}

\paragraph{Детермінована стартова конфігурація (матчинг публікацій).}
Таблиця~\ref{tab:features-sim-weights} надає відтворювану початкову конфігурацію ваг як відправну точку (не як емпіричне твердження). Порогові значення нижче мають бути налаштовані.

\paragraph{Порогові значення рішень (потрібна калібровка).}
Використовуємо два пороги $\theta_{\mathrm{match}} > \theta_{\mathrm{review}}$:
\begin{equation}
  \textsf{auto-merge if } m(r_i,r_j)\ge \theta_{\mathrm{match}},\qquad
  \textsf{flag if } \theta_{\mathrm{review}}\le m(r_i,r_j) < \theta_{\mathrm{match}}.
\end{equation}
Типові стартові значення: $\theta_{\mathrm{match}}=0.90$, $\theta_{\mathrm{review}}=0.70$ (\emph{підлягає емпіричній калібровці}).

\subsection{Кластеризація}
\label{sec:gr-clustering}

Побудуємо ваговий неорієнтований граф $G=(V,E)$, де $V=R$, а $E=\{(r_i,r_j, m_{ij}) : (r_i,r_j)\in\mathcal{P}\}$ з вагою ребра $m_{ij}=m(r_i,r_j)$.

\paragraph{Базова кластеризація через порогові компоненти зв’язності.}
Визначимо $E_{\theta}=\{(i,j): m_{ij}\ge\theta_{\mathrm{cluster}}\}$ та візьмемо зв’язні компоненти $G_{\theta}=(V,E_{\theta})$ як кластери. Типове стартове значення: $\theta_{\mathrm{cluster}}=\theta_{\mathrm{match}}$.

\paragraph{Обмеження проти ланцюгових злиттів (рекомендовано).}
Чисте транзитивне замикання може спричиняти \emph{ланцюгові злиття} (A близьке до B, B близьке до C, але A далеко від C). Рекомендуємо забезпечити принаймні одне з таких обмежень:
\begin{enumerate}
  \item \textbf{Уточнення complete-linkage:} зливати дві компоненти лише якщо $\min\{m(u,v): u\in c_a, v\in c_b\}\ge\theta_{\mathrm{cluster}}$.
  \item \textbf{Узгодженість із центром:} підтримувати тимчасове центроїдне представлення $\hat{g}(c)$ (наприклад, нормалізована назва та рік) і вимагати $m(r,\hat{g}(c))\ge\theta_{\mathrm{centroid}}$ для всіх елементів.
\end{enumerate}
Ці параметри ($\theta_{\mathrm{centroid}}$ тощо) \emph{підлягають налаштуванню}.

\subsection{Злиття атрибутів (побудова Golden Record)}
\label{sec:gr-fusion}

Для кластера $c=\{r_1,\dots,r_{|c|}\}$ Golden Record є відображенням атрибутів
\begin{equation}
  g(c) = \big\{ (a, \hat{v}_a(c)) : a\in\mathcal{A}_\tau \big\},
\end{equation}
де $\hat{v}_a(c)$ обчислюється атрибут-специфічним оператором злиття.

\paragraph{Надійність джерел і довіра до поля.}
Нехай $\rho(s,a)\in[0,1]$ позначає надійність джерела $s$ для атрибута $a$ (визначається політикою; може бути сталою для джерела або навченою). Нехай $q(r,a)\in[0,1]$ --- сигнал якості значення на рівні запису (наприклад, наявність, валідність формату, актуальність за часовим штампом). Визначимо ефективну вагу голосу:
\begin{equation}
  w(r,a) = \rho(s(r),a)\cdot q(r,a).
\end{equation}

\paragraph{Оператори злиття.}
Для поширених типів атрибутів:
\begin{itemize}
  \item \textbf{Ідентифікатори (DOI, ORCID, ROR):} якщо валідний ідентифікатор існує, обрати за максимумом ваги; якщо спостерігаються кілька конфліктних валідних ID, розв’язати конфлікт через правило розв’язання нічиїх (нижче) та позначити конфлікт.
  \item \textbf{Рядки (назва, venue):} обрати значення, що максимізує зважений бал узгодженості
  \begin{equation}
    \hat{v}_a(c) = \arg\max_{v\in\mathcal{V}_a(c)} \sum_{r\in c:\,x_r[a]=v} w(r,a),
  \end{equation}
  з опційною пост-нормалізацією (відновлення капіталізації).
  \item \textbf{Числові значення (кількість цитувань):} застосувати консервативний агрегатор, наприклад $\max$ по джерелах, фіксуючи всі значення джерел у провенансі.
  \item \textbf{Множини (ключові слова):} об’єднання з канонікалізацією та опційним обрізанням за частотою.
  \item \textbf{Списки (автори):} розв’язувати як задачу вирівнювання послідовностей (за наявності) або брати об’єднання з упорядкуванням за більшістю; неоднозначні випадки мають переходити на ручну перевірку.
\end{itemize}

\subsection{Розв’язання конфліктів і політика розв’язання нічиїх}
\label{sec:gr-tiebreak}

Визначимо детерміновану політику розв’язання нічиїх $\mathrm{TB}(a)$ для кожного атрибута $a$. Для метаданих публікацій приклад політики (має бути документований і версіонований) такий:
\begin{equation}
  \text{DOI наявний і валідний} \succ \text{реєстри на кшталт Crossref} \succ \text{курована CRIS} \succ \text{OpenAlex} \succ \text{самозвіт ORCID}.
\end{equation}

\paragraph{Жорсткі обмеження.}
Щоб запобігти катастрофічним злиттям, накладаємо жорсткі обмеження (fail-fast):
\begin{itemize}
  \item Якщо обидва записи мають валідний DOI і вони різні, тоді \emph{не виконувати auto-merge} без явного зовнішнього встановлення еквівалентності DOI (рідкісні випадки; зазвичай потрібна ручна перевірка).
  \item Якщо рік відрізняється більш ніж на $\Delta_y$ (типово 2), вимагати ручної перевірки, якщо DOI не збігається.
\end{itemize}

\subsection{Провенанс (на рівні атрибутів) та відображення на W3C PROV}
\label{sec:gr-provenance}

\paragraph{Схема провенансу на рівні атрибутів.}
Для кожного злитого атрибута $a$ у Golden Record $g(c)$ зберігаємо мінімальний кортеж провенансу:
\begin{equation}
  \mathrm{prov}_a(c) = \Big( \hat{v}_a,\; \{(s,\mathrm{id}_s, v, t)\}_{\text{support}},\; \mathrm{rule}_a,\; \mathrm{ver}(\Theta),\; \gamma_a \Big),
\end{equation}
де \emph{support} перелічує сирі спостереження, що підтримують значення, $\mathrm{rule}_a$ --- ідентифікатор правила злиття (включно з tie-break), $\mathrm{ver}(\Theta)$ --- геш версії конфігурації, а $\gamma_a\in[0,1]$ --- необов’язковий бал упевненості (визначений правилом або моделлю).

\paragraph{Відображення на W3C PROV.}
Відобразимо конвеєр на W3C PROV через такі відповідності:
\begin{itemize}
  \item \textbf{prov:Entity}
  \begin{itemize}
    \item кожен сирий запис $r$ (як інжестований),
    \item кожен проміжний нормалізований запис $\tilde{r}$,
    \item кожен Golden Record $g(c)$, а також, опційно, кожне значення атрибута $\hat{v}_a$ як окрема сутність.
  \end{itemize}
  \item \textbf{prov:Activity}
  \begin{itemize}
    \item нормалізація $\mathsf{Normalize}$,
    \item blocking $\mathsf{Block}$,
    \item скоринг $\mathsf{ScorePair}$,
    \item кластеризація $\mathsf{Cluster}$,
    \item злиття $\mathsf{Fuse}$.
  \end{itemize}
  \item \textbf{prov:Agent}
  \begin{itemize}
    \item джерела даних $s\in\mathcal{S}$ (як організації/сервіси),
    \item реалізація конвеєра (software agent) з версією,
    \item опційний людський куратор (якщо наявні HITL-рішення).
  \end{itemize}
\end{itemize}

Як мінімум, стверджуємо:
\begin{align}
  &\texttt{prov:wasGeneratedBy}(g(c), \mathsf{Fuse}),\\
  &\texttt{prov:used}(\mathsf{Fuse}, \tilde{r})\ \forall r\in c,\\
  &\texttt{prov:wasAttributedTo}(\tilde{r}, s(r)),\\
  &\texttt{prov:wasDerivedFrom}(g(c), r)\ \forall r\in c.
\end{align}
Для трасування на рівні атрибутів представляємо кожне $\hat{v}_a$ як сутність із посиланнями \texttt{prov:wasDerivedFrom} на підтримувальні спостереження.

\subsection{Крайові випадки та операційні запобіжники}
\label{sec:gr-edgecases}

Наведені нижче випадки мають бути явно (детерміновано) оброблені та залоговані для аудиту:
\begin{enumerate}
  \item \textbf{Відсутні або невалідні ідентифікатори:} перейти на ознаки назва/рік/автори; підвищити частоту ручної перевірки.
  \item \textbf{Препринт проти опублікованої версії:} трактувати як (a) різні сутності, пов’язані через \texttt{isVersionOf}, або (b) одну сутність із версіонованими атрибутами; обрати політику та дотримуватися її послідовно.
  \item \textbf{Переклади назв / багатомовні варіанти:} дозволити поля альтернативних назв; застосовувати нормалізацію з урахуванням транслітерації; уникати агресивного stemming, якщо він погіршує міжмовний матчинг.
  \item \textbf{Омонімія авторів:} вимагати додаткових доказів (афіліація, граф співаторів, ORCID) перед злиттям персон.
  \item \textbf{Варіанти назв venue та зміни серій:} підтримувати таблиці псевдонімів; логувати правила канонікалізації.
  \item \textbf{Ризик ланцюгового злиття:} увімкнути анти-ланцюгові обмеження (розділ~\ref{sec:gr-clustering}) та ручну перевірку для прикордонних випадків.
  \item \textbf{Конфлікти між високонадійними джерелами:} якщо два джерела з високою надійністю не узгоджуються щодо ключового атрибута (наприклад, DOI), позначити як конфлікт і не виконувати авто-розв’язання.
\end{enumerate}

\subsection{Таблиця «ознака $\to$ подібність $\to$ вага» (стартова конфігурація)}
\label{sec:gr-feature-table}

Таблиця~\ref{tab:features-sim-weights} задає детермінований базовий варіант, який можна реалізувати без навчання. Значення є \emph{стартовими} та мають бути емпірично калібровані для цільового корпусу.

\begin{table}[t]
\centering
\small
\begin{tabular}{p{0.26\linewidth} p{0.30\linewidth} p{0.10\linewidth} p{0.28\linewidth}}
\toprule
\textbf{Ознака} & \textbf{Подібність / предикат} & \textbf{Вага} & \textbf{Обґрунтування}\\
\midrule
Збіг DOI & $\mathbb{I}[\mathrm{doi}_i=\mathrm{doi}_j]$ (після нормалізації) & 3.0 & Сильний глобальний ідентифікатор; майже вирішальний за валідності.\\
Конфлікт DOI & $\mathbb{I}[\mathrm{doi}_i\neq\mathrm{doi}_j \wedge \text{обидва валідні}]$ & $-5.0$ & Жорсткий негативний доказ; запобігає катастрофічним злиттям.\\
Косинус назви (TF--IDF) & $\cos(\mathrm{tfidf}(t_i),\mathrm{tfidf}(t_j))$ & 1.2 & Стійко до малих правок; сильний сигнал за відсутності DOI.\\
Символьна подібність назви & $1-\mathrm{normLevenshtein}(t_i,t_j)$ & 0.8 & Уловлює друкарські помилки та різницю в пунктуації.\\
Узгодженість року & $\mathbb{I}[|y_i-y_j|\le 1]$ & 0.6 & Допомагає розрізненню; толерантно до зсувів індексації.\\
Узгодженість venue & Jaccard по токенах venue & 0.4 & Корисно, але шумно через скорочення.\\
Перетин авторів & Jaccard по нормалізованих прізвищах & 0.9 & Критично, коли назви короткі/неоднозначні.\\
Узгодженість сторінок & $\mathbb{I}[p_i=p_j]$ (якщо наявно) & 0.2 & Слабкий підтримувальний доказ; часто відсутній.\\
Збіг URL/домену & точний домен або рівність нормалізованого URL & 0.2 & Корисно для препринтів та OA-джерел.\\
\bottomrule
\end{tabular}
\caption{Детермінований базовий варіант для зв’язування записів публікацій (стартові ваги; налаштувати емпірично).}
\label{tab:features-sim-weights}
\end{table}

\subsection{Відтворювана конфігурація (гіперпараметри)}
\label{sec:gr-hyperparams}

Для відтворюваності конвеєр має публікувати конфігураційний бандл (YAML/JSON), що містить щонайменше:
\begin{itemize}
  \item параметри нормалізації $\Theta_{\nu,\tau}$ (токенізація, стоп-слова, regex для DOI, політика Unicode),
  \item параметри blocking (визначення $\mathcal{B}$; $\Delta_y$; minhash seeds; top-$K$, якщо використовується retrieval),
  \item параметри скорингу ($\beta$ або геш чекпойнта моделі; визначення ознак),
  \item пороги ($\theta_{\mathrm{match}}, \theta_{\mathrm{review}}, \theta_{\mathrm{cluster}}$),
  \item політику злиття (надійності джерел $\rho(s,a)$; порядок розв’язання нічиїх $\mathrm{TB}(a)$; агрегатори для кожного атрибута),
  \item версію схеми провенансу та налаштування експорту PROV.
\end{itemize}

Усі значення параметрів, що потребують емпіричного налаштування, мають бути явно позначені в рукописі як такі, що \emph{підлягають налаштуванню на розмічених development-даних}, та не подані як фінальні результати.
